\begin{tikzpicture}[
    >=Latex,
    font=\small,
    nota/.style={font=\footnotesize, align=center},
    flujo/.style={-{Latex[width=5pt,length=4pt]}, thick, gray},
    celda/.style={draw, thick, fill=gray!30}
]

% Tolva genérica: centro de la salida en (#1,#2)
\newcommand{\tolvap}[2]{%
    \draw[thick] (#1-0.9,#2+1.3) -- (#1+0.9,#2+1.3) -- (#1+0.25,#2+0.3) --
        (#1+0.25,#2) (#1-0.25,#2) -- (#1-0.25,#2+0.3) -- (#1-0.9,#2+1.3);
}

% (a) Volumétrico: masa = densidad x volumen
\tolvap{0}{2.2}
\draw[thick] (-0.25,1.9) rectangle (1.3,2.2);
\foreach \x in {-0.1,0.25,0.6,0.95}{\draw (\x,1.9) -- (\x+0.2,2.2);}
\draw[flujo] (1.15,1.85) -- (1.15,1.2);
\draw[thick] (0.75,0.55) rectangle (1.55,1.15);
\node[nota] at (0.4,3.85) {Vueltas o tiempo};
\node[nota] at (0.5,0.05) {(a) Volumétrico};

% (b) Pesaje por lotes: el recipiente sobre la celda gana peso
\tolvap{4.2}{2.2}
\draw[thick] (3.95,1.9) rectangle (5.5,2.2);
\foreach \x in {4.1,4.45,4.8,5.15}{\draw (\x,1.9) -- (\x+0.2,2.2);}
\draw[flujo] (5.35,1.85) -- (5.35,1.2);
\draw[thick] (4.95,0.75) rectangle (5.75,1.15);
\draw[celda] (5.05,0.5) rectangle (5.65,0.75);
\node[nota] at (4.6,3.85) {Mide la masa que entra};
\node[nota] at (4.7,0.05) {(b) Pesaje por lotes};

% (c) Pérdida de peso: la tolva sobre celdas pierde peso
\tolvap{8.6}{2.2}
\draw[celda] (7.55,3.5) rectangle (7.85,3.75);
\draw[celda] (9.35,3.5) rectangle (9.65,3.75);
\draw[thick] (8.35,1.9) rectangle (9.9,2.2);
\foreach \x in {8.5,8.85,9.2,9.55}{\draw (\x,1.9) -- (\x+0.2,2.2);}
\draw[flujo] (9.75,1.85) -- (9.75,1.2);
\node[nota] at (8.65,4.05) {Mide la masa que sale};
\node[nota] at (9.1,0.05) {(c) Pérdida de peso};

\end{tikzpicture}
